\section{练习}

\begin{question}\begin{drill}[有限直导线的长线误差]\label{mg1:drill:1}
长为 $2l$ 的直导线通有电流 $I$，场点在其中垂线上，垂距为 $a$。若用无限长导线公式近似，并以无限长值为分母定义相对误差，要求该误差不超过5\%，至少要求 $l/a$ 多大？
\end{drill}
\fig{fig:mg1:drill:1:s}{题设：有限直导线的长线误差}{\begin{tikzpicture}[scale=.88]
\draw[world](0,-2)--(0,2);\fill(3,0)circle(2pt)node[right]{$P$};\draw[dim](0,0)--(3,0)node[midway,below]{$a$};\node[left]at(0,2){$l$};\node[left]at(0,-2){$-l$};\draw[flow](0,.6)--(0,1.5)node[left]{$I$};\draw[guide](0,2)--(3,0)--(0,-2);
\end{tikzpicture}}
\end{question}
\answerlink{mg1:drill:1}
\begin{question}\begin{drill}[圆弧与径向引线的圆心场]\label{mg1:drill:2}
一段半径 $R$、圆心角 $3\pi/2$ 的圆弧通有逆时针电流 $I$，两端经径向引线引至远处闭合，忽略远处回路的场。求圆心场。
\end{drill}
\fig{fig:mg1:drill:2:s}{题设：圆弧与径向引线的圆心场}{\begin{tikzpicture}[scale=.88]
\draw[world](0:1.5)arc(0:270:1.5);\draw[world](1.5,0)--(3,0) (0,-1.5)--(0,-3);\draw[flow](45:1.5)arc(45:100:1.5);\fill(0,0)circle(2pt)node[right]{$O$};
\end{tikzpicture}}
\end{question}
\answerlink{mg1:drill:2}
\begin{question}\begin{drill}[正多边形线圈的圆心场]\label{mg1:drill:3}
正 $N$ 边形线圈外接圆半径为 $R$，电流为 $I$。求圆心场，并检验 $N\to\infty$ 的极限。
\end{drill}
\fig{fig:mg1:drill:3:s}{题设：正多边形线圈的圆心场}{\begin{tikzpicture}[scale=.88]
\foreach\a in{0,60,...,300}{\draw[world](\a:1.6)--({\a+60}:1.6);}\draw[dim](0,0)--(1.6,0)node[midway,above]{$R$};\node[below]at(0,0){$O$};
\end{tikzpicture}}
\end{question}
\answerlink{mg1:drill:3}
\begin{question}\begin{drill}[旋转带电圆盘的中心场]\label{mg1:drill:4}
半径 $R$ 的绝缘薄圆盘面电荷密度为 $\sigma>0$，以角速度 $\omega$ 绕轴逆时针匀速转动。求轴上磁场的积分式及圆心值。
\end{drill}
\fig{fig:mg1:drill:4:s}{题设：旋转带电圆盘的中心场}{\begin{tikzpicture}[scale=.88]
\draw[world](0,0)circle(1.6);\fill(0,0)circle(1.4pt)node[below]{$O$};\draw[dim](0,0)--(1.6,0)node[midway,below]{$R$};\draw[flow](1.6*.707,1.6*.707)arc(45:115:1.6);\fill[blue!15,even odd rule](0,0)circle(.98)(0,0)circle(.85);\draw[guide](0,0)circle(.85)(0,0)circle(.98);\draw[flow](-2,1.2)--(-.7,.63)node[pos=0,above left]{窄环 $r,\dd r$};\node at(0,2.2){$\sigma,\omega$};
\end{tikzpicture}}
\end{question}
\answerlink{mg1:drill:4}
\begin{question}\begin{drill}[中空圆柱电流的三段场]\label{mg1:drill:5}
无限长圆柱导体占据 $a<r<b$，轴向电流密度均为 $J$。求空心、导体内及外部的磁场，并检查两个边界。
\end{drill}
\fig{fig:mg1:drill:5:s}{题设：中空圆柱电流的三段场}{\begin{tikzpicture}[scale=.88]
\draw[world](0,0)circle(1.6);\fill(0,0)circle(1.4pt)node[below]{$O$};\draw[dim](0,0)--(1.6,0)node[midway,below]{$b$};\draw[flow](1.6*.707,1.6*.707)arc(45:115:1.6);\draw[world](0,0)circle(.7);\node at(.3,.4){$a$};\node at(-1,0.7){$\odot J$};
\end{tikzpicture}}
\end{question}
\answerlink{mg1:drill:5}
\begin{question}\begin{drill}[厚电流层的连续磁场]\label{mg1:drill:6}
沿 $x,z$ 方向无限延展、厚度为 $2a$ 的平板占据 $-a<y<a$，均匀电流密度 $\bm J=J\bm e_z$。外无其他电流或外加场。求所有位置的 $\bm B(y)$。
\end{drill}
\fig{fig:mg1:drill:6:s}{题设：厚电流层的连续磁场}{\begin{tikzpicture}[scale=.88]
\fill[gray!15](-2,-1)rectangle(2,1);\draw[world](-2,-1)--(2,-1) (-2,1)--(2,1);\node at(0,0){$\odot J$};\node[right]at(2,1){$a$};\node[right]at(2,-1){$y=-a$};\draw[guide](-2.3,0)--(2.3,0);\draw[axis](-2.7,-1.7)--(-2.7,1.8)node[above]{$y$};\draw[axis](-2.4,-1.5)--(.1,-1.5)node[right]{$x$};
\end{tikzpicture}}
\end{question}
\answerlink{mg1:drill:6}
\begin{question}\begin{drill}[同轴正反电流的空间磁场]\label{mg1:drill:7}
半径 $a$ 的实心圆柱通有均匀电流 $I$ 向外，半径 $b>a$ 的极薄同轴圆柱面上通有总电流 $I$ 向里。求三段场，说明外场为零的原因。
\end{drill}
\fig{fig:mg1:drill:7:s}{题设：同轴正反电流的空间磁场}{\begin{tikzpicture}[scale=.88]
\draw[world](0,0)circle(1.7);\fill(0,0)circle(1.4pt)node[below]{$O$};\draw[dim](0,0)--(1.7,0)node[midway,below]{$b$};\draw[flow](1.7*.707,1.7*.707)arc(45:115:1.7);\draw[world](0,0)circle(.7);\node at(0,.4){$\odot I$};\node at(2.5,1.3){外层 $\otimes I$};\draw[dim](0,0)--(-.606,.35)node[midway,above left]{$a$};
\end{tikzpicture}}
\end{question}
\answerlink{mg1:drill:7}
\begin{question}\begin{drill}[离散线排的近场修正]\label{mg1:drill:8}
细线间距为 $d$，各沿 $+z$ 通电流 $I$，横截面为 $(kd,0)$，$k$ 为整数。求 $(0,y)$、$y>0$ 处的场。可用恒等式 $\sum_{k=-\infty}^{\infty}[(kd)^2+y^2]^{-1}=\pi\coth(\pi y/d)/(dy)$。高度多大时，离散场相对连续面场的增加不超过1\%？
\end{drill}
\fig{fig:mg1:drill:8:s}{题设：离散线排的近场修正}{\begin{tikzpicture}[scale=.88]
\draw[world](-2.5,0)--(2.5,0);\foreach\x in{-2,-1,0,1,2}{\draw[world](\x,0)circle(.1);\fill(\x,0)circle(1pt);}\fill(0,2)circle(2pt)node[right]{$P(0,y)$};\node at(0,-.6){间距 $d$};
\end{tikzpicture}}
\end{question}
\answerlink{mg1:drill:8}
